% Generated draft: requires mathematical and source review.
\documentclass[11pt]{article}
\usepackage[margin=25mm]{geometry}
\usepackage{fontspec}
\setmainfont{Latin Modern Roman}
\usepackage{amsmath,amssymb,mathtools,unicode-math}
\setmathfont{Latin Modern Math}
\usepackage{enumitem,needspace,xurl,hyperref}
\setlength{\parindent}{0pt}
\setlength{\parskip}{5pt}
\setlength{\emergencystretch}{3em}
\newcommand{\N}{\mathbb{N}}
\newcommand{\Z}{\mathbb{Z}}
\newcommand{\Q}{\mathbb{Q}}
\newcommand{\R}{\mathbb{R}}
\newcommand{\C}{\mathbb{C}}
\newcommand{\F}{\mathbb{F}}
\newcommand{\A}{\mathbb{A}}
\newcommand{\PP}{\mathbb{P}}
\newcommand{\E}{\mathbb{E}}
\newcommand{\eps}{\varepsilon}
\newcommand{\dd}{\,\mathrm{d}}
\begin{document}
\section*{UnsolvedMath ID: 11100030}
\textbf{Source problem code:} catalogue entry
\subsection{Definitions and mathematical statement}
Let $G$ be a finite $p$-group and let $k$ be a field of characteristic $p$, the intended modular setting. The excerpt does not explicitly state this characteristic assumption. Write $\operatorname{stmod}(kG)$ for finitely generated, equivalently finite-dimensional, $kG$-modules with morphisms taken modulo maps factoring through projectives. The group algebra is self-injective; in this $p$-group modular setting its finitely generated projective modules are free. The finite-module category must not be confused with the larger stable category of all modules.

A thick subcategory is a full triangulated subcategory closed under direct summands. Put $R=H^*(G;k)$ and let $\operatorname{Proj}R$ be its homogeneous projective prime spectrum. A subset is specialization closed when it contains every specialization of each of its points; it need not itself be one closed subset. Using the standard cohomological action, set
\[
 V_G(M)=\{\mathfrak p\in\operatorname{Proj}R:
 \operatorname{Ext}_{kG}^*(M,M)_{\mathfrak p}\ne0\}.
\]

\textbf{Applications problem 6 (Hovey).} Improve the Benson--Carlson--Rickard classification, in particular by removing algebraic closedness of $k$. The requested classification for arbitrary fields in the modular setting is the correspondence
\[
 \mathcal C\longmapsto\bigcup_{M\in\mathcal C}V_G(M),\qquad
 Y\longmapsto\{M\in\operatorname{stmod}(kG):V_G(M)\subseteq Y\},
\]
between thick subcategories and specialization-closed subsets $Y\subseteq\operatorname{Proj}H^*(G;k)$. The source recalls the algebraically closed case as the Benson--Carlson--Rickard theorem and proposes removing that restriction, possibly using stable-category Galois theory. It supplies no separate formal Galois-theoretic conjecture.

The distinction between thick and tensor-ideal thick subcategories is material. For a modular $p$-group, the unique simple module is $k$, so tensoring a fixed object with a finite module built from $k$ preserves its thick closure. This explains the usual all-thick formulation here. For a general finite group, a support classification stated for tensor ideals must not be promoted silently to all thick subcategories, and projective modules need not be free.

The registry reports a solved-in-literature extension through Benson--Iyengar--Krause, but gives no exact theorem or paper and phrases it broadly for finite groups over arbitrary fields. That broader claim requires these tensor-ideal, characteristic, and small-versus-big category qualifications. The original $p$-group field-extension task remains the precise source request, with solved status retained as an attributed report.
\subsection{Short English statement}
Extend the support classification of thick finite-module stable subcategories for finite $p$-groups from algebraically closed fields to arbitrary modular fields. Preserve specialization-closed supports and distinguish the tensor-ideal statement needed for general finite groups.
\subsection{Sources}
\begin{itemize}
\item Local UnsolvedMath registry, permanent ID 11100030, source code catalogue entry; attributed to Mark Hovey, no proposal year supplied.
\item \emph{Hovey's Algebraic Topology Problem List}, Applications, item 6; surviving mirror \url{https://www-users.cse.umn.edu/~tlawson/hovey/}. The registry records HTML-ordered-list extraction and access on 2026-07-29; the mirror has 65 top-level items on nine category pages.
\item The source names Benson--Carlson--Rickard and Hovey's work with Palmieri and Strickland, without full bibliography entries. The research summary attributes the reported extension to Benson--Iyengar--Krause, also without a precise citation.
\end{itemize}
\end{document}
